\section*{अध्याय \ref{ch:g9:maps-of-the-earth} --- पूरी पृथ्वी के मानचित्र}

\begin{solution}{exo:g9:maps-of-the-earth:1}
(क) \omterm{def:g9:maps-of-the-earth:projection}{मानचित्र प्रक्षेप}। (ख) \omterm{def:g9:maps-of-the-earth:contour}{समोच्च रेखा}। (ग) \omterm{def:g9:maps-of-the-earth:numerical-scale}{संख्यात्मक मापनी}। (घ) \omterm{def:g9:maps-of-the-earth:contour}{समोच्च
अंतराल}। (ङ) \omterm{def:g9:maps-of-the-earth:remote-sensing}{सुदूर संवेदन}।
\end{solution}

\begin{solution}{exo:g9:maps-of-the-earth:2}
(क) $4 \times 25\,000 = \num{100000}$~cm $= \qty{1}{km}$। (ख) $3 \times
50\,000 = \num{150000}$~cm $= \qty{1.5}{km}$। (ग) $2.5 \times 1\,000\,000
= \num{2500000}$~cm $= \qty{25}{km}$। (घ) $\qty{1}{km} = \num{100000}$~cm;
$100\,000 \div 25\,000 = \qty{4}{cm}$।
\end{solution}

\begin{solution}{exo:g9:maps-of-the-earth:3}
(क) ग़लत: किसी गोले को लंबाइयाँ विकृत किए बिना चपटा नहीं किया जा सकता, इसलिए हर \omterm{def:g3:continents-and-oceans:world-map}{विश्व
मानचित्र} कुछ न कुछ विकृत करता है। (ख) सही। (ग) ग़लत: समक्षेत्र \omterm{def:g2:drawing-to-map:map}{मानचित्र} क्षेत्रफल बनाए रखता
है, पर आकारों को मोड़ देता है, ख़ासकर अपने किनारों के पास। (घ) सही।
\end{solution}

\begin{solution}{exo:g9:maps-of-the-earth:4}
(क) ढाल खड़ी है। (ख) ज़मीन हल्की ढाल वाली है। (ग) एक \omterm{def:g1:landscapes-around-us:landform}{पहाड़ी}, जिसकी चोटी सबसे छोटी रेखा के
अंदर है। (घ) ज़मीन के किसी बिंदु की केवल एक ऊँचाई होती है, इसलिए वह अलग-अलग ऊँचाइयों वाली
दो रेखाओं पर नहीं हो सकता।
\end{solution}

\begin{solution}{exo:g9:maps-of-the-earth:5}
(क) \ang{60} उत्तर पर \omterm{def:g1:landscapes-around-us:water}{झील} 4 गुना ज़्यादा बड़ी बनाई जाती है: वह
$4 \times 1000 = \qty{4000}{km^2}$ जितनी बड़ी दिखती है। \omterm{def:g3:continents-and-oceans:equator}{भूमध्य रेखा} पर क्षेत्रफल बढ़ाए नहीं
जाते, इसलिए उसी आकार में बनी \omterm{def:g1:landscapes-around-us:water}{झील} का सच्चा क्षेत्रफल \qty{4000}{km^2} है।
(ख) $4 \div 1.33 \approx 3$ गुना बड़ा। (ग) \omterm{def:g3:continents-and-oceans:poles}{ध्रुवों} की ओर क्षेत्रफल, और उत्तर से दक्षिण की
दूरियाँ, बिना सीमा के बढ़ते जाते हैं: \omterm{def:g3:continents-and-oceans:poles}{ध्रुव} \omterm{def:g3:continents-and-oceans:equator}{भूमध्य रेखा} से अनंत दूरी पर
होते।
\end{solution}

\begin{solution}{exo:g9:maps-of-the-earth:6}
(क) \qty{50}{m}। (ख) शिखर \qty{544}{m} पर है; उसके चारों ओर की सूचक \omterm{def:g9:maps-of-the-earth:contour}{समोच्च रेखा}
\qty{500}{m} की रेखा है। (ग) पश्चिम में: \omterm{def:g1:landscapes-around-us:water}{नदी} ढाल पर नीचे बहती है, इसलिए उसका तल धारा
की दिशा में आगे, पूर्व की ओर, नीचा होता है। (घ) लगभग 300 और \qty{400}{m} की रेखाओं के
बीच: वहाँ ढाल सबसे खड़ी है, \omterm{def:g1:landscapes-around-us:landform}{घाटी} के तल के पास भी और शिखर के पास भी ज़्यादा
हल्की।
\end{solution}

\begin{solution}{exo:g9:maps-of-the-earth:7}
(क) C लगभग \qty{400}{m} पर, D लगभग \qty{530}{m} पर। (ख) लगभग
\qty{135}{m}, C से लगभग \qty{1.3}{km} पर, \omterm{def:g1:landscapes-around-us:landform}{घाटी} के तल पर। (ग) लगभग
0.3 से \qty{0.4}{km}। (घ) दूरियों की तुलना में ऊँचाइयाँ छोटी हैं: बिना खींचे \omterm{def:g1:landscapes-around-us:landform}{पहाड़ियाँ} लगभग
समतल दिखतीं।
\end{solution}

\begin{solution}{exo:g9:maps-of-the-earth:8}
(क) $2 \times 25\,000 = \num{50000}$~cm $= \qty{500}{m}$। (ख) $300 - 200 =
\qty{100}{m}$। (ग) $100 \div 500 = 0.2$, \qty{20}{\%} की ढाल। (घ) ज़्यादा हल्का:
\qty{4}{cm} का अर्थ \qty{1000}{m} है, इसलिए ढाल $100 \div 1000 = \qty{10}{\%}$ है,
आधी खड़ी।
\end{solution}

\begin{solution}{exo:g9:maps-of-the-earth:9}
(क) $14 \times 18 = 252$ परिक्रमाएँ। (ख) $14 \times 103 = 1442$ मिनट, लगभग
\qty{24}{h}: एक दिन। (ग) $5.5 \div 1 = 5.5$ गुना। (घ) उदाहरण के लिए: पीछे हटता
\omterm{def:g8:ice-on-earth:glacier}{हिमनद}, कटा हुआ जंगल, बढ़ता कस्बा, खिसकता \omterm{def:g1:landscapes-around-us:water}{तट} या \omterm{def:g1:landscapes-around-us:water}{नदी}।
\end{solution}

\begin{solution}{exo:g9:maps-of-the-earth:10}
(क) $20\,200 \div 300\,000 \approx 0.067$~s $= 6.7 \times 10^{-2}$~s।
(ख) $300\,000 \times 10^{-3} = \qty{300}{km}$। (ग) समय की एक नन्ही त्रुटि दूरी
की बहुत बड़ी त्रुटि बनाती है: घड़ियों को अत्यंत सटीक होना पड़ता है।
\end{solution}

\begin{solution}{exo:g9:maps-of-the-earth:11}
(क) बिंदु A तक की दूरी: $\sqrt{3.6^2 + 4.8^2} = \sqrt{12.96 + 23.04} =
\sqrt{36} = 6$; बिंदु B तक की दूरी: $\sqrt{6.4^2 + 4.8^2} = \sqrt{40.96 +
23.04} = \sqrt{64} = 8$। (ख) $(3.6, -4.8)$, रेखा AB के पार दर्पण-प्रतिबिंब।
(ग) C से $(3.6, 4.8)$ तक दूरी: $\sqrt{3.6^2 + 5.2^2} =
\sqrt{40} \approx 6.3$; $(3.6, -4.8)$ तक: $\sqrt{3.6^2 + 14.8^2} =
\sqrt{232} \approx 15.2$। अभिग्राही $(3.6, 4.8)$ पर है। (घ) एक दूरी
एक वृत्त देती है (अंतरिक्ष में एक गोला), दो दूरियाँ दो संभव बिंदु देती हैं; चुनने के लिए एक और
दूरी की ज़रूरत होती है।
\end{solution}

\begin{solution}{exo:g9:maps-of-the-earth:12}
(क) $29 \div 2.1 \approx 14$ गुना। (ख) $2.1 \times 10.5 \approx
22$~मिलियन~\unit{km^2}। (ग) ग्रीनलैंड लगभग 22~मिलियन
\unit{km^2} घेरता दिखता है, अफ़्रीका के 29~मिलियन के क़रीब: मर्केटर दीवार \omterm{def:g2:drawing-to-map:map}{मानचित्र} पर
दोनों लगभग एक ही आकार के दिखते हैं, हालाँकि अफ़्रीका 14 गुना बड़ा है। (घ) एक
\omterm{def:g9:maps-of-the-earth:projection}{समक्षेत्र प्रक्षेप}।
\end{solution}

\begin{solution}{pb:g9:maps-of-the-earth:1}
\textbf{1.} $50\,000$~cm $= \qty{500}{m}$।

\textbf{2.} $3.6 \times 500 = \qty{1800}{m} = \qty{1.8}{km}$।

\textbf{3.} B--W: $2.5 \times 500 = \qty{1250}{m}$; W--T: $2.35 \times 500
\approx \qty{1180}{m}$।

\textbf{4.} उत्तर-पूर्व की ओर।

\textbf{5.} \qty{50}{m}।

\textbf{6.} S--B नीचे जाता है: $296 - 143 = \qty{153}{m}$।

\textbf{7.} B--W: $335 - 143 = \qty{192}{m}$; W--T: $544 - 335 =
\qty{209}{m}$।

\textbf{8.} $209 \div 1180 \approx 0.18$, लगभग \qty{18}{\%}।

\textbf{9.} S से T तक की सीधी रेखा \omterm{def:g1:landscapes-around-us:water}{नदी} को वहाँ पार करती जहाँ कोई पुल नहीं है; रास्ता B पर
पुल तक नीचे जाता है।

\textbf{10.} $1.8 + 1.25 + 1.18 \approx \qty{4.2}{km}$, \qty{4}{km} प्रति घंटे की चाल पर: लगभग 63 मिनट; चढ़ाई $192 + 209 = \qty{401}{m}$: लगभग 40 मिनट
और। कुल मिलाकर लगभग 1~घंटा~43~मिनट।

\textbf{11.} T से वह B तक नीचे जाती है, फिर B से S तक चढ़ती है: $296 -
143 = \qty{153}{m}$।

\textbf{12.} लंबाई: $1.8 + 1.25 + 1.18 \approx \textbf{4.2 km}$; कुल
चढ़ाई: $192 + 209 = \textbf{401 m}$।
\end{solution}
